\subsection{How much stronger is the prover?}\label{sec:separations}

\paragraph{A universal verifier.} Fix a consistent theory $T$ (say ZFC). Let a
$V_T$-argument be a randomized program $A$, a time bound $t$, a budget
$\delta_A\le\delta$, and a $T$-proof of the bounded arithmetic sentence
``at most $\delta_A2^t$ coin strings make $A^{(t)}(C,D)$ output an interval
missing $\sum_xD(x)C(x)$''; $V_T$ checks the proof and runs $A$. It is sound
whenever $T$ is (a false bounded sentence has a true $\Sigma_1$ negation,
which $T$ would refute), and it $p$-simulates every argument system whose
soundness $T$ proves, by instantiating that proof. So it can use any FPRAS,
compiler, mixing-time bound or combinatorial theorem as a certified rule. This
is the analogue, for estimation, of proof systems that are optimal relative to
a theory \cite{kp1989}.

\paragraph{Lower bounds: the prover helps unless $\NP\subseteq\BPP$.}
Uniquely satisfiable circuits have $p=2^{-n}$, certified exactly by one
witness. A sound Task-1 procedure that outputs any $L>0$ on all of them
decides SAT via Valiant--Vazirani \cite{vv1986}. Under one-way permutations
(with a Goldreich--Levin hard-core bit \cite{gl1989}),
$C_y(x,r)=[(\pi(x),r,\langle x,r\rangle)=y]$ gives an average-case version
with $(1+\eps)$-tight certificates.

\paragraph{Upper bounds: the prover helps exactly when $P$ is not weakly
automatable.} For a refutation system $P$ let $Z_P^m$ be the circuits $C$ with
a $P$-refutation of size $m(|C|)$. A sound Task-1 procedure that outputs
$[0,0]$ on all of $Z_P^m$ exists iff the canonical pair $(Z_P^m,\mathrm{SAT})$ is
separable in randomized polynomial time (if it outputs $[0,0]$, the circuit
cannot be satisfiable, by soundness; conversely run the separator); separability of
the canonical pair is weak automatability of $P$ \cite{pudlak2003pairs}. Padding
turns this into a separation for positive, exactly certified $p$: with
$\mathrm{pad}_a(C)(\beta,x,w)=(\beta\wedge x=0^n\wedge w=0^a)\vee(\neg\beta\wedge
C(x))$, $p=2^{-n-a-1}$ exactly when $C$ is unsatisfiable (certified by one point
and $P$'s refutation), while $p\ge2^{-n-1}$ when $C$ is satisfiable, so any Task-1
procedure complete within a factor $2^a$ separates the canonical pair. Known
results then give:

\begin{corollary}\label{cor:automatability}
If factoring Blum integers is hard for polynomial-size circuits, then for
$P\in\{\mathrm{TC}^0\text{-Frege},\text{Frege}\}$ \cite{bpr2000}, and for
extended Frege if RSA is secure \cite{kp1998}, there are circuit families on
which the PCIS verifier with $P$ certifies $p$ exactly while no sound
polynomial-time procedure without a prover is complete, even within a factor
$2^a$. Under LWE the same holds against quantum procedures for
$\mathrm{TC}^0$-Frege \cite{acg2024}.
\end{corollary}

Finding the short proofs is the hard part: even for resolution, finding a
refutation of size polynomial in the shortest one is $\NP$-hard
\cite{atserias2020automating}. And for preimage circuits $[G(x)=y]$ of
pseudorandom generators, short proofs that $y\notin\mathrm{range}(G)$ are
conjectured not to exist in strong proof systems (proof complexity generators
\cite{abrw2004,krajicek2004dual}), so there the verifier cannot beat sampling for
upper bounds either.

\paragraph{Two-sided, exact.} Let $L\in\mathsf{UP}\cap\coNP$ with verifiers $A$
(unique witness, $0^k$ never a witness) and $B$, and
$C_u(\beta,v)=(\beta\wedge v=0^k)\vee(\neg\beta\wedge A(u,v))$, so
$p(C_u)=2^{-k-1}(1+[u\in L])$. With $T$ proving uniqueness and disjointness,
$V_T$ certifies both values exactly (by the witness of $A$ or of $B$), while a
Task-1 procedure complete on $\{C_u\}$ decides $L$. For bits of the
prime factorisation (AKS-certified primes) this gives a two-sided separation
under ``factoring $\notin\BPP$''.

\paragraph{Why not NP-hardness?} If SAT reduces (randomised, many-one) to the
gap problem ``$p\ge\gamma t$ versus $p\le t$'' with unsatisfiable formulas
mapped to the \emph{small} side, and some MA system is complete for upper
bounds on that side, then $\coNP\subseteq\AM$: Arthur samples the reduction,
Merlin sends arguments, Arthur accepts if most of them certify $U<\gamma t$.
With the opposite orientation, Goldwasser--Sipser already puts the problem
in $\AM$. Either way $\PH$ collapses \cite{bhz1987}. So upper-bound and
two-sided separations must come from structured, non-$\NP$-hard assumptions
(factoring, RSA, LWE), as above, while lower-bound separations may come from
$\NP\not\subseteq\BPP$.

\paragraph{Universality.} Upper bounds: complete for all circuits only if
$\coNP\subseteq\MA$ (\cref{thm:ma}). Lower bounds: complete for all circuits
within a constant factor iff $\mathrm{prSBP}\subseteq\mathrm{prMA}$
\cite{bgm2006}; this fails relative to an oracle, and holds under derandomisation
hypotheses strong enough to fool nondeterministic circuits, with a verifier
whose soundness then depends on the hypothesis.

